\documentclass[12pt,a4paper]{article}

% ===== PACKAGES =====
\usepackage[utf8]{inputenc}
\usepackage[vietnamese]{babel}
\usepackage{amsmath, amssymb, amsfonts}
\usepackage{geometry}
\usepackage{enumitem}
\usepackage{array}
\usepackage{booktabs}
\usepackage{fancyhdr}
\usepackage{titlesec}
\usepackage{xcolor}
\usepackage{tcolorbox}

\geometry{a4paper, margin=2.5cm}

% ===== CUSTOM STYLES =====
\definecolor{mainblue}{RGB}{0,90,160}
\definecolor{darkgreen}{RGB}{0,120,60}
\definecolor{resultbg}{RGB}{230,245,255}

\newtcolorbox{resultbox}{
  colback=resultbg, colframe=mainblue,
  boxrule=0.8pt, arc=3pt,
  left=8pt, right=8pt, top=6pt, bottom=6pt,
  fonttitle=\bfseries, title=Kết quả
}

\newtcolorbox{notebox}[1][]{
  colback=yellow!8, colframe=orange!70!black,
  boxrule=0.6pt, arc=3pt,
  left=8pt, right=8pt, top=6pt, bottom=6pt,
  fonttitle=\bfseries, title={Lưu ý}, #1
}

\pagestyle{fancy}
\fancyhf{}
\fancyhead[L]{\textit{Phân tích \& Thiết kế Thuật toán}}
\fancyhead[R]{\textit{GV: Huỳnh Thị Thanh Thương}}
\fancyfoot[C]{\thepage}

\title{\textbf{Bài giải Bài 2: Xác định $O$, $\Theta$, $\Omega$ và Phân lớp hàm}\\[6pt]
\large (Có lập luận/giải thích theo hướng dẫn giảng viên)}
\author{}
\date{}

\begin{document}
\maketitle
\thispagestyle{fancy}

% =========================================
% CÂU 1: Đúng / Sai
% =========================================
\section{Câu 1: Xác định đúng hay sai}

\textbf{Đề bài:} Dùng định nghĩa không hình thức (informal) của $O$, $\Theta$, $\Omega$ để xác định các khẳng định sau đúng hay sai.

\bigskip

Trước hết, ta rút gọn hàm cần xét:
\[
  \frac{n(n+1)}{2} = \frac{n^2 + n}{2} \in \Theta(n^2).
\]
Bỏ hệ số $1/2$, bỏ bậc thấp $n$, phần chính là $n^2$. Vậy bậc tăng trưởng là $n^2$.

\bigskip

\begin{enumerate}[label=\textbf{\alph*.}]

  \item $\dfrac{n(n+1)}{2} \in O(n^3)$ \hfill \textcolor{darkgreen}{\textbf{ĐÚNG}} \\[4pt]
  \textbf{Lập luận:} Ta đã biết $\frac{n(n+1)}{2} \in \Theta(n^2)$.
  Mà $n^2 \prec n^3$ (bậc $2 < 3$), nên $n^2 \in O(n^3)$.

  Theo tính chất bắc cầu: $\frac{n(n+1)}{2} \in O(n^2) \subseteq O(n^3)$.

  (Như cô giảng trong buổi 4: một hàm có thể thuộc nhiều lớp $O$ --- $O(n^2)$, $O(n^3)$, $O(n^4)$,\ldots\ đều đúng.
  Hàm bậc 2 bị chặn trên bởi $n^3$ là hiển nhiên.)

  \bigskip

  \item $\dfrac{n(n+1)}{2} \in O(n^2)$ \hfill \textcolor{darkgreen}{\textbf{ĐÚNG}} \\[4pt]
  \textbf{Lập luận:} $\frac{n(n+1)}{2} = \frac{n^2 + n}{2} \leq \frac{n^2 + n^2}{2} = n^2$ với mọi $n \geq 1$.

  Chọn $c = 1$, $n_0 = 1$: ta có $\frac{n(n+1)}{2} \leq 1 \cdot n^2$ với mọi $n \geq 1$.

  $\Longrightarrow$ Theo định nghĩa $O$: $\frac{n(n+1)}{2} \in O(n^2)$. \checkmark

  \bigskip

  \item $\dfrac{n(n+1)}{2} \in \Theta(n^3)$ \hfill \textcolor{red}{\textbf{SAI}} \\[4pt]
  \textbf{Lập luận (phản chứng --- theo cô hướng dẫn buổi 5):}

  Giả sử khẳng định đúng, tức $\frac{n(n+1)}{2} \in \Theta(n^3)$.

  Theo định nghĩa $\Theta$: $\exists\, c_1, c_2 > 0,\; n_0 \in \mathbb{N}$ sao cho:
  \[
    c_1 \cdot n^3 \leq \frac{n(n+1)}{2} \leq c_2 \cdot n^3 \quad \forall\, n \geq n_0.
  \]

  Xét vế trái, chia cả hai vế cho $n^2$ (với $n \geq 1$):
  \[
    c_1 \cdot n \leq \frac{n+1}{2}.
  \]

  Bất đẳng thức $c_1 \cdot n^3 \leq \frac{n^2+n}{2}$ suy ra:
  \[
    c_1 \leq \frac{n^2+n}{2n^3} = \frac{1}{2n} + \frac{1}{2n^2} \to 0 \quad \text{khi } n \to \infty.
  \]

  $\Longrightarrow$ Không tồn tại $c_1 > 0$ thỏa mãn với \textbf{mọi} $n \geq n_0$. \textbf{Mâu thuẫn!}

  Vậy giả sử sai $\Rightarrow \frac{n(n+1)}{2} \notin \Theta(n^3)$.

  \bigskip

  \item $\dfrac{n(n+1)}{2} \in \Omega(n)$ \hfill \textcolor{darkgreen}{\textbf{ĐÚNG}} \\[4pt]
  \textbf{Lập luận:} $\frac{n(n+1)}{2} = \frac{n^2+n}{2} \geq \frac{n}{2}$ với mọi $n \geq 0$.

  Chọn $c = \frac{1}{2}$, $n_0 = 0$: ta có $\frac{n(n+1)}{2} \geq \frac{1}{2} \cdot n$ với mọi $n \geq 0$.

  $\Longrightarrow$ Theo định nghĩa $\Omega$: $\frac{n(n+1)}{2} \in \Omega(n)$. \checkmark

  (Bậc tăng trưởng $n^2$ lớn hơn $n$, nên hiển nhiên $n^2$ là chặn dưới của $\Omega(n)$.)

\end{enumerate}

\begin{resultbox}
\renewcommand{\arraystretch}{1.4}
\begin{tabular}{llcl}
  \textbf{a.} & $n(n+1)/2 \in O(n^3)$      & $\Longrightarrow$ & \textcolor{darkgreen}{\textbf{ĐÚNG}} (bậc 2 $<$ bậc 3) \\
  \textbf{b.} & $n(n+1)/2 \in O(n^2)$      & $\Longrightarrow$ & \textcolor{darkgreen}{\textbf{ĐÚNG}} (cùng bậc, chặn trên chặt) \\
  \textbf{c.} & $n(n+1)/2 \in \Theta(n^3)$ & $\Longrightarrow$ & \textcolor{red}{\textbf{SAI}} (phản chứng: $c_1 \to 0$) \\
  \textbf{d.} & $n(n+1)/2 \in \Omega(n)$   & $\Longrightarrow$ & \textcolor{darkgreen}{\textbf{ĐÚNG}} (bậc 2 $>$ bậc 1) \\
\end{tabular}
\end{resultbox}

% =========================================
% CÂU 2: Phân lớp Θ
% =========================================
\newpage
\section{Câu 2: Phân lớp $\Theta(g(n))$ cho từng hàm}

\textbf{Đề bài:} Cho biết mỗi hàm sau thuộc lớp $\Theta(g(n))$ nào (chọn $g(n)$ đơn giản nhất). Chứng minh.

\bigskip

\begin{enumerate}[label=\textbf{\alph*.}]

  % ----- (a) -----
  \item $(n^2 + 1)^{10}$

  \textbf{Phân lớp:} Khai triển theo bậc cao nhất:
  \[
    (n^2 + 1)^{10} = n^{20} + \underbrace{10n^{18} + \cdots + 1}_{\text{bậc thấp hơn}}.
  \]
  Bậc cao nhất theo $n$ là $n^{20}$. Bỏ hệ số, bỏ bậc thấp.

  $\Longrightarrow \boxed{(n^2+1)^{10} \in \Theta(n^{20})}$

  \textbf{Chứng minh bằng định nghĩa:}

  \textit{Chặn trên} ($O$): $(n^2+1)^{10} \leq (n^2 + n^2)^{10} = (2n^2)^{10} = 2^{10} \cdot n^{20} = 1024\,n^{20}$,
  với mọi $n \geq 1$.

  Chọn $c_2 = 1024$, $n_0 = 1$.

  \textit{Chặn dưới} ($\Omega$): $(n^2+1)^{10} \geq (n^2)^{10} = n^{20}$, với mọi $n \geq 0$.

  Chọn $c_1 = 1$, $n_0 = 0$.

  Vậy $c_1 n^{20} \leq (n^2+1)^{10} \leq c_2 n^{20}$ với $c_1=1$, $c_2=1024$, $n_0=1$. \checkmark

  \bigskip

  % ----- (b) -----
  \item $\sqrt{10n^2 + 7n + 3}$

  \textbf{Phân lớp:} Bên trong căn, bậc cao nhất là $10n^2$.
  \[
    \sqrt{10n^2 + 7n + 3} \approx \sqrt{10n^2} = \sqrt{10}\,n.
  \]
  Bỏ hệ số $\sqrt{10}$: bậc tăng trưởng là $n$.

  $\Longrightarrow \boxed{\sqrt{10n^2 + 7n + 3} \in \Theta(n)}$

  \textbf{Chứng minh:}

  \textit{Chặn trên:} Với $n \geq 1$: $7n \leq 7n^2$ và $3 \leq 3n^2$. Vậy:
  \[
    10n^2 + 7n + 3 \leq 10n^2 + 7n^2 + 3n^2 = 20n^2.
  \]
  \[
    \sqrt{10n^2 + 7n + 3} \leq \sqrt{20}\,n \approx 4{,}47\,n.
  \]
  Chọn $c_2 = 5$, $n_0 = 1$.

  \textit{Chặn dưới:} $10n^2 + 7n + 3 \geq 10n^2$ (vì $7n + 3 > 0$ khi $n \geq 0$). Vậy:
  \[
    \sqrt{10n^2 + 7n + 3} \geq \sqrt{10}\,n \approx 3{,}16\,n.
  \]
  Chọn $c_1 = 3$, $n_0 = 1$.

  Vậy $3n \leq \sqrt{10n^2+7n+3} \leq 5n$ với mọi $n \geq 1$. \checkmark

  \bigskip

  % ----- (c) -----
  \item $2n\lg(n+2)^2 + (n+2)^2 \lg\dfrac{n}{2}$

  \textbf{Phân lớp:} Rút gọn từng phần:

  \textit{Phần 1:} $2n\lg(n+2)^2 = 2n \cdot 2\lg(n+2) = 4n\lg(n+2)$.

  Khi $n$ đủ lớn: $\lg(n+2) \approx \lg n$. Vậy Phần 1 $\in \Theta(n\log n)$.

  \textit{Phần 2:} $(n+2)^2\lg\frac{n}{2} = (n+2)^2(\lg n - \lg 2) = (n+2)^2(\lg n - 1)$.

  Khi $n$ đủ lớn: $(n+2)^2 \approx n^2$ và $\lg n - 1 \approx \lg n$. Vậy Phần 2 $\in \Theta(n^2\log n)$.

  \textit{Gộp lại (quy tắc cộng --- lấy max):}
  \[
    \max\bigl(\Theta(n\log n),\; \Theta(n^2\log n)\bigr) = \Theta(n^2\log n).
  \]

  $\Longrightarrow \boxed{2n\lg(n+2)^2 + (n+2)^2\lg\frac{n}{2} \in \Theta(n^2\lg n)}$

  \textbf{Chứng minh (sketch):}

  \textit{Chặn trên:} Với $n \geq 2$:
  \begin{align*}
    &4n\lg(n+2) + (n+2)^2(\lg n - 1) \\
    &\leq 4n \cdot 2\lg n + (2n)^2 \cdot \lg n \quad (\text{vì } n+2 \leq 2n,\; \lg(n+2) \leq 2\lg n) \\
    &= 8n\lg n + 4n^2\lg n = (8n + 4n^2)\lg n \leq 5n^2\lg n.
  \end{align*}
  Chọn $c_2 = 5$, $n_0$ đủ lớn.

  \textit{Chặn dưới:} Chỉ cần lấy Phần 2:
  \[
    (n+2)^2(\lg n - 1) \geq n^2 \cdot \frac{1}{2}\lg n = \frac{1}{2}n^2\lg n, \quad \text{với } n \geq 4.
  \]
  Chọn $c_1 = \frac{1}{2}$, $n_0 = 4$. \checkmark

  \bigskip

  % ----- (d) -----
  \item $2^{n+1} + 3^{n-1}$

  \textbf{Phân lớp --- đưa về cùng cơ số 2:}

  \begin{itemize}
    \item $2^{n+1} = 2 \cdot 2^n$. Phần mũ: $n$, hệ số: $2$. Thuộc $\Theta(2^n)$.
    \item $3^{n-1} = \frac{1}{3} \cdot 3^n = \frac{1}{3} \cdot 2^{n\lg 3} \approx \frac{1}{3} \cdot 2^{1{,}585\,n}$. Phần mũ: $1{,}585\,n$. Thuộc $\Theta(3^n)$.
  \end{itemize}

  So phần mũ: $n < 1{,}585\,n$ nên $2^n \prec 3^n$.

  Quy tắc cộng (lấy max): $2^{n+1} + 3^{n-1} \in \Theta(3^n)$.

  $\Longrightarrow \boxed{2^{n+1} + 3^{n-1} \in \Theta(3^n)}$

  \textbf{Chứng minh:}

  \textit{Chặn trên:} Với $n \geq 1$: $2^{n+1} \leq 3^n$ (vì $2^{n+1} = 2 \cdot 2^n \leq 2 \cdot 3^n$). Vậy:
  \[
    2^{n+1} + 3^{n-1} \leq 2 \cdot 3^n + \frac{1}{3} \cdot 3^n = \frac{7}{3} \cdot 3^n.
  \]
  Chọn $c_2 = \frac{7}{3}$, $n_0 = 1$.

  \textit{Chặn dưới:} $2^{n+1} + 3^{n-1} \geq 3^{n-1} = \frac{1}{3} \cdot 3^n$.

  Chọn $c_1 = \frac{1}{3}$, $n_0 = 1$.

  Vậy $\frac{1}{3}\cdot 3^n \leq 2^{n+1} + 3^{n-1} \leq \frac{7}{3}\cdot 3^n$ với mọi $n \geq 1$. \checkmark

  \bigskip

  % ----- (e) -----
  \item $\lfloor \log_2 n \rfloor$

  \textbf{Phân lớp:} Hàm phần nguyên (floor) của $\log_2 n$.

  Ta biết: $\log_2 n - 1 < \lfloor \log_2 n \rfloor \leq \log_2 n$, với mọi $n \geq 1$.

  Vậy $\lfloor \log_2 n \rfloor$ xấp xỉ $\log_2 n$ (chỉ sai lệch tối đa 1 đơn vị --- không đáng kể khi $n$ lớn).

  $\Longrightarrow \boxed{\lfloor \log_2 n \rfloor \in \Theta(\lg n)}$

  \textbf{Chứng minh:}

  \textit{Chặn trên:} $\lfloor \log_2 n \rfloor \leq \log_2 n = 1 \cdot \lg n$. Chọn $c_2 = 1$, $n_0 = 1$.

  \textit{Chặn dưới:} $\lfloor \log_2 n \rfloor > \log_2 n - 1 \geq \frac{1}{2}\log_2 n$, với mọi $n \geq 4$ (vì $\lg n \geq 2$ khi $n \geq 4$).

  Chọn $c_1 = \frac{1}{2}$, $n_0 = 4$. \checkmark

\end{enumerate}

\begin{resultbox}
\renewcommand{\arraystretch}{1.5}
\begin{tabular}{lll}
  \textbf{a.} & $(n^2+1)^{10}$      & $\in \Theta(n^{20})$ \\
  \textbf{b.} & $\sqrt{10n^2+7n+3}$ & $\in \Theta(n)$ \\
  \textbf{c.} & $2n\lg(n+2)^2 + (n+2)^2\lg\frac{n}{2}$ & $\in \Theta(n^2\lg n)$ \\
  \textbf{d.} & $2^{n+1} + 3^{n-1}$ & $\in \Theta(3^n)$ \\
  \textbf{e.} & $\lfloor\log_2 n\rfloor$ & $\in \Theta(\lg n)$ \\
\end{tabular}
\end{resultbox}

% =========================================
% CÂU 3: Chứng minh đa thức bậc k ∈ Θ(n^k)
% =========================================
\newpage
\section{Câu 3: Chứng minh đa thức bậc $k$ thuộc $\Theta(n^k)$}

\textbf{Đề bài:} Chứng minh rằng mọi đa thức bậc $k$:
\[
  p(n) = a_k n^k + a_{k-1} n^{k-1} + \cdots + a_1 n + a_0, \quad \text{với } a_k > 0,
\]
đều thuộc $\Theta(n^k)$.

\bigskip

\textbf{Cần chứng minh:} Tồn tại $c_1, c_2 > 0$ và $n_0 \in \mathbb{N}$ sao cho:
\[
  c_1 \cdot n^k \leq p(n) \leq c_2 \cdot n^k, \quad \forall\, n \geq n_0.
\]

\subsection*{Chứng minh chặn trên: $p(n) \in O(n^k)$}

Với mọi $n \geq 1$, ta có $n^i \leq n^k$ khi $0 \leq i \leq k$. Do đó:
\begin{align*}
  p(n) &= a_k n^k + a_{k-1}n^{k-1} + \cdots + a_1 n + a_0 \\
       &\leq |a_k| n^k + |a_{k-1}| n^k + \cdots + |a_1| n^k + |a_0| n^k \\
       &= \Bigl(\sum_{i=0}^{k} |a_i|\Bigr) \cdot n^k.
\end{align*}

Chọn $c_2 = \displaystyle\sum_{i=0}^{k} |a_i|$, $n_0 = 1$.

$\Longrightarrow p(n) \leq c_2 \cdot n^k$ với mọi $n \geq 1$. \checkmark

\subsection*{Chứng minh chặn dưới: $p(n) \in \Omega(n^k)$}

Chia cả hai vế cho $n^k$ (với $n \geq 1$):
\[
  \frac{p(n)}{n^k} = a_k + \frac{a_{k-1}}{n} + \frac{a_{k-2}}{n^2} + \cdots + \frac{a_0}{n^k}.
\]

Khi $n \to \infty$, các phần $\frac{a_i}{n^{k-i}} \to 0$ với mọi $i < k$. Nên:
\[
  \lim_{n\to\infty} \frac{p(n)}{n^k} = a_k > 0.
\]

Do giới hạn bằng $a_k > 0$, tồn tại $n_0$ đủ lớn sao cho:
\[
  \frac{p(n)}{n^k} \geq \frac{a_k}{2} > 0, \quad \forall\, n \geq n_0.
\]

Suy ra:
\[
  p(n) \geq \frac{a_k}{2} \cdot n^k, \quad \forall\, n \geq n_0.
\]

Chọn $c_1 = \dfrac{a_k}{2}$.

$\Longrightarrow p(n) \geq c_1 \cdot n^k$ với mọi $n \geq n_0$. \checkmark

\subsection*{Kết luận}

\begin{resultbox}
Từ hai phần trên, ta có:
\[
  \frac{a_k}{2} \cdot n^k \;\leq\; p(n) \;\leq\; \Bigl(\sum_{i=0}^{k}|a_i|\Bigr) \cdot n^k, \quad \forall\, n \geq n_0.
\]
Với $c_1 = \dfrac{a_k}{2} > 0$, \; $c_2 = \displaystyle\sum_{i=0}^{k}|a_i| > 0$, \; $n_0$ đủ lớn.

Theo định nghĩa $\Theta$:
\[
  \boxed{p(n) = a_k n^k + a_{k-1}n^{k-1} + \cdots + a_0 \;\in\; \Theta(n^k).}
\]
\end{resultbox}

\begin{notebox}[title={Ý nghĩa thực tiễn (theo cô Thương)}]
Kết quả này xác nhận phương pháp ``ước lượng nhanh'' mà cô đã dạy:
\begin{itemize}
  \item Khi gặp một đa thức, chỉ cần nhìn \textbf{bậc cao nhất} theo $n$.
  \item Bỏ hết hệ số và các bậc thấp hơn.
  \item Kết luận ngay: $\Theta(n^k)$ --- với $k$ là bậc cao nhất.
\end{itemize}
Ví dụ: $0.001n^4 + 3n^3 + 1 \in \Theta(n^4)$ --- không cần biến đổi thêm!
\end{notebox}

\end{document}
